% ==============================================================
\section{結論}\label{sec:conclusion}
% ==============================================================

本研究では、頻出パターンのサポート変動と外部イベントの帰属を
空間次元に拡張する\textbf{時空間イベント帰属パイプライン}を提案した。
サポート曲面上の位置別変化点検出、
時間的・空間的近接度を統合した帰属スコア、
および空間構造を保存する置換検定により、
「どの地域で・いつ・なぜパターン頻度が変わったか」に答える枠組みを提供する。

合成データ実験では、十分な信号強度（$\beta \geq 0.3$）において
1次元パイプラインに対しF1を最大+0.80改善し（$\beta = 0.5$で$F1 = 0.89$）、
Precisionで一貫した優位性を示した。

\paragraph{制約と限界.}
(1) 空間次元はグリッド離散化を前提としており、
連続空間の解像度選択が結果に影響する。
(2) 位置あたりの観測数が少ない場合（$\beta \leq 0.2$）、
変化点の統計的有意性が低下しRecallが低くなる（\ref{sec:discussion}節）。
(3) 次元数 $k$ の増加に伴い計算量が $O(2^{k+1})$ で増大し、
$k \geq 4$ では実用性が制限される。

\paragraph{今後の課題.}
(1) 位置あたりの信号が薄い場合の適応的空間粒度選択。
(2) イベントの空間スコープが未知の場合の自動推定。
(3) ストリーミング設定へのオンライン拡張。
(4) 空間スコープのランダム化を含む、
より一般的な帰無仮説の設計。
